\textbf{- 参考资料：}
\href{https://mysite.science.uottawa.ca/mnevins/papers/StephenHarrigan2017LWE.pdf}{格密码与
带误差学习问题}~\cite{lattice-crypto}


$ $

格密码在明确的计算困难假设及参数选择下，为后量子安全提供候选方案；不能无条件地声称它能够抵抗所有量子攻击。本节介绍后续使用的 LWE（带误差学习）、RLWE（环带误差学习）、GLWE、GLev 和 GGSW 构造。





\subsection{概述}
\label{subsec:lattice-overview}

本概述中的所有等式均在 $\mathbb Z_q$ 上成立。掩码独立均匀采样；误差先从指定分布采样为小整数，仅在运算时模 $q$ 约减。

假设有一个未知的 $k$ 维向量 $\vec{s}$ 作为秘密密钥，以及许多公开已知的 $k$ 维向量 $\vec{a}^{\langle i \rangle}$。

再假设有大量如下点积 $\vec{s} \cdot \vec{a}^{\langle i \rangle}$：

$\vec{s} \cdot \vec{a}^{\langle 0 \rangle} = s_0\cdot a_{0}^{\langle 0 \rangle} + s_1\cdot a_{1}^{\langle 0 \rangle} + \cdots + s_{k-1}\cdot a_{k-1}^{\langle 0 \rangle} = b^{\langle 0 \rangle}$

$\vec{s} \cdot \vec{a}^{\langle 1 \rangle} = s_0\cdot a_{0}^{\langle 1 \rangle} + s_1\cdot a_{1}^{\langle 1 \rangle} + \cdots + s_{k-1}\cdot a_{k-1}^{\langle 1 \rangle} = b^{\langle 1 \rangle}$

$\vec{s} \cdot \vec{a}^{\langle 2 \rangle} = s_0\cdot a_{0}^{\langle 2 \rangle} + s_1\cdot a_{1}^{\langle 2 \rangle} + \cdots + s_{k-1}\cdot a_{k-1}^{\langle 2 \rangle} = b^{\langle 2 \rangle}$

\text{ } $\vdots$

$ $

在无噪声情形下，若 $k$ 个样本的系数矩阵在 $\mathbb Z_q$ 上可逆，它们就能确定 $\vec s$。当 $q$ 为素数时，这要求矩阵满秩；当模数为合数时，行列式必须为单位。仅有 $k$ 个方程，并不能保证方程独立或唯一恢复。

但是，如果在上述每个方程中都随机加入一个未知的小噪声 $e^{\langle i \rangle}$，即误差，得到：

$b^{\langle 0 \rangle} = s_0\cdot a_{0}^{\langle 0 \rangle} + s_1\cdot a_{1}^{\langle 0 \rangle} + \cdots + s_{k-1}\cdot a_{k-1}^{\langle 0 \rangle} + e^{\langle 0 \rangle} \pmod q$

$b^{\langle 1 \rangle} = s_0\cdot a_{0}^{\langle 1 \rangle} + s_1\cdot a_{1}^{\langle 1 \rangle} + \cdots + s_{k-1}\cdot a_{k-1}^{\langle 1 \rangle} + e^{\langle 1 \rangle} \pmod q$

$b^{\langle 2 \rangle} = s_0\cdot a_{0}^{\langle 2 \rangle} + s_1\cdot a_{1}^{\langle 2 \rangle} + \cdots + s_{k-1}\cdot a_{k-1}^{\langle 2 \rangle} + e^{\langle 2 \rangle} \pmod q$

\text{ } $\vdots$


$ $

噪声使得这些方程不能直接作为关于 $\vec s$ 的精确线性方程求解。若 $k$ 个误差各有 $r$ 种可能值，枚举全部 $r^k$ 个误差向量是一种可行攻击。这只是该枚举方法的工作量上界，并非所有算法的下界，也不是安全性证明。额外样本可能帮助攻击者，而且攻击者始终可以忽略它们。LWE 困难性是针对明确维数、模数、秘密与误差分布以及样本数量的计算假设；过小或具有特殊结构的误差可能使某些实例易解。

基于这一直观认识，格密码所依赖的数学困难问题可描述如下：

$ $

\begin{tcolorbox}[title={\textbf{\tboxlabel{\ref*{subsec:lattice-overview}} LWE（带误差学习）与 RLWE 问题}}]
\textbf{\underline{LWE 问题}}

考虑如下形式的样本：$b = \vec{s} \cdot\vec{a} + e$，其中 $e$ 是后面将介绍的小噪声。

对每个样本，从 $\mathbb Z_q^k$ 中均匀选取 $\vec a$，并独立地从明确指定的分布 $\chi$ 采样小整数误差 $e$，随后将 $e$ 模 $q$ 约减。秘密 $\vec s$ 从其指定分布采样一次并重复使用。误差并非在 $\mathbb Z_q$ 上均匀分布。安全实验允许获得关于安全参数的多项式数量的这类样本。

$(\vec{a}^{\langle 1 \rangle}, b^{\langle 1 \rangle})$，其中 $b^{\langle 1 \rangle} = \vec{s} \cdot \vec{a}^{\langle 1 \rangle} + e^{\langle 1 \rangle}$

$(\vec{a}^{\langle 2 \rangle}, b^{\langle 2 \rangle})$，其中 $b^{\langle 2 \rangle} = \vec{s} \cdot \vec{a}^{\langle 2 \rangle} + e^{\langle 2 \rangle}$

$(\vec{a}^{\langle 3 \rangle}, b^{\langle 3 \rangle})$，其中 $b^{\langle 3 \rangle} = \vec{s} \cdot \vec{a}^{\langle 3 \rangle} + e^{\langle 3 \rangle}$

\text{ } $\vdots$

假设攻击者拥有大量二元组 $(\vec{a}^{\langle i \rangle}, b^{\langle i \rangle})$。在这一设定下，以下困难问题构成 LWE（带误差学习）密码体制的安全基础：

$ $

\begin{itemize}
\item \textbf{搜索问题：}给定这些样本，恢复秘密 $\vec s$。搜索 LWE 假设认为，对于所选参数族，不存在以不可忽略概率成功的高效算法。
\item \textbf{判定问题：}区分 LWE 样本 $(\vec a,\langle\vec a,\vec s\rangle+e)$ 与 $\mathbb Z_q^k\times\mathbb Z_q$ 中独立均匀的二元组。判定 LWE 假设认为，对于所选参数族，每个高效区分器的优势都可忽略。
\end{itemize}

$ $

搜索到判定的归约需要前提；这两种表述并不对所有模数和分布无条件等价。例如，Regev 的引理 4.2 要求模数为关于维数多项式有界的素数，并满足其中的误差及预言机条件。其最坏情形格问题归约也有明确的噪声和模数要求，而且是量子归约。参见 \href{https://cims.nyu.edu/~regev/papers/qcrypto.pdf}{Regev, On Lattices, Learning with Errors, Random Linear Codes, and Cryptography} 的定理 1.1 及第 4 节。

$ $

\textbf{\underline{RLWE 问题}}

在 RLWE 的代数形式中，使用 $\mathbb Z_q[X]/(X^n+1)$ 上的 $B=AS+E$，多项式代表元的次数小于 $n$。搜索问题要求恢复 $S$，判定问题要求将这些二元组与均匀环元素对区分开。仅替换记号并不能转移所有 LWE 安全性定理：环、模数以及误差和秘密分布必须满足相应环 LWE 结论的前提。

\end{tcolorbox}




\noindent\textbf{小秘密的适用范围：}正规形归约可以把均匀 LWE 秘密改成从\emph{误差分布}中采样的秘密。对于素数幂模数，仅根据掩码选择可逆矩阵 $A_0$，对应 $\vec b_0=A_0\vec s+\vec e_0$。对新鲜行向量样本 $(\vec a,b)$，令
\[
\vec a'=-\vec aA_0^{-1},\qquad b'=b-\vec aA_0^{-1}\vec b_0
=\vec a'\vec e_0+e.
\]
新的掩码仍均匀，新秘密的分布为 $\chi^k$；均匀输入对也映射成均匀对。预处理材料是此归约的内部状态。在环版本中，同一恒等式要求初始环掩码可逆。参见 \href{https://web.cs.ucla.edu/~sahai/work/web/2009\%20Publications/ACPS09.pdf}{Applebaum--Cash--Peikert--Sahai，第 3.1 节引理 2}。该论证不会把高斯误差变成均匀三元秘密。因此，后续三元秘密构造使用针对实际分布的判定假设，除非另给分布匹配的归约。维数计数或“小”这个形容词都不能代替该归约。

\subsection{LWE 密码体制}
\label{subsec:lattice-scheme}

下面的正确性公式假设 $t\mid q$，取 $\Delta=q/t$，并以 $[x]_q$ 表示 $x\bmod q$ 的中心化整数代表元。舍入作用于该代表元，解码结果还需模 $t$ 约减。所有位移描述仅适用于 $\Delta$ 为二幂的情形；一般的缩放是乘以或除以 $\Delta$。完整密文是 $(\vec a,b)$，并非只有 $b$。

LWE 密码体制使用加密公式 $b = \vec{s} \cdot \vec{a} + \Delta \cdot m + e$，其中 $\vec{s}$ 为秘密密钥，$\vec{a}$ 为每次加密时选取的公开随机向量，$m$ 为明文，$e$ 为每次加密时从正态分布中随机选取的小噪声，$b$ 为密文。$\Delta$ 是明文 $M$ 的尺度因子，相当于将 $m$ 向左移动 $\text{log}_2\Delta$ 个二进制位。加密前，先将明文左移若干位，即 $\text{log}_2\Delta$ 位，为低位中的误差预留足够空间。

\begin{figure}[h!]
    \centering
  \includegraphics[width=0.8\linewidth]{figures/TFHE-fig1.pdf}
  \caption{LWE 中的明文缩放与加噪示意图：$\Delta \cdot m + e \in \mathbb{Z}_q$}
  \label{fig:scaling}
\end{figure}

\autoref{fig:scaling} 直观展示了 $\Delta \cdot m + e$：将明文 $m$ 左移 $\text{log}_2\Delta$ 位，并加入噪声 $e$。实际的加密和解密公式如下：

\begin{tcolorbox}[title={\textbf{\tboxlabel{\ref*{subsec:lattice-scheme}} 基于格的 LWE 密码体制}}]
\begin{itemize}
    \item \textbf{\underline{加密}：}$b^{\langle i \rangle} = \vec{s} \cdot \vec{a}^{\langle i \rangle} + \Delta \cdot m^{\langle i \rangle} + e^{\langle i \rangle}$，其中 $b^{\langle i \rangle}$ 和 $\vec{a}^{\langle i \rangle}$ 对攻击者也是公开已知的，而 $\vec{s}, m^{\langle i \rangle}, e^{\langle i \rangle}$ 则未知，仅秘密密钥持有者知晓。

    $ $
    
    \item \textbf{\underline{解密}：}$\left\lceil [b^{\langle i\rangle}-\vec s\cdot\vec a^{\langle i\rangle}]_q/\Delta\right\rfloor\bmod t=m^{\langle i\rangle}$ $\Big($，前提是 $|e^{\langle i \rangle}| < \dfrac{\Delta}{2}\Big)$。
\end{itemize}
\end{tcolorbox}

$\lfloor \rceil_\Delta$ 表示将数舍入到最接近的 $\Delta$ 的整数倍。例如，$\lfloor 16 \rceil_{10} = 20$ 表示将 16 舍入到最接近的 10 的整数倍。再如，$\lfloor 17 \rceil_{8} = 16$ 表示将 17 舍入到最接近的 8 的整数倍；由于 17 更接近 16 而非 24，因此结果为 16。


$ $

\noindent \textbf{正确性：}在解密过程中，计算 $b^{\langle i \rangle} - \vec{s} \cdot \vec{a}^{\langle i \rangle}$ 得到 $\Delta \cdot m^{\langle i \rangle}+ e^{\langle i \rangle}$，如 \autoref{fig:scaling} 所示。随后计算 $\lceil \Delta \cdot m^{\langle i \rangle} + e^{\langle i \rangle} \rfloor_{\Delta}$，即将数值舍入到最接近的 $\Delta$ 的整数倍；只要所加噪声满足 $|e^{\langle i \rangle}| < \frac{\Delta}{2}$，就会得到 $\Delta \cdot m^{\langle i \rangle}$。也就是说，如果噪声小于 $\frac{\Delta}{2}$，它就会在舍入过程中被消去。最后，将 $\Delta \cdot m^{\langle i \rangle}$ 右移 $\text{log}_2 \Delta$ 位，即可得到 $m^{\langle i \rangle}$。因此，只要保证 $|e^{\langle i \rangle}| < \frac{\Delta}{2}$，也就是噪声 $e^{\langle i \rangle}$ 必须小于相应阈值，就能够在解密的舍入过程中消去 $e^{\langle i \rangle}$，恢复原始的 $\Delta \cdot m^{\langle i \rangle}$。将 $m^{\langle i \rangle}$ 按 $\Delta$ 缩放有两个目的：(i) 加密时在低位为 $e^{\langle i \rangle}$ 留出空间，使噪声位与明文位不相互干扰，避免破坏明文；(ii) 解密时消去低位中的噪声 $e^{\langle i \rangle}$，同时保持明文 $m^{\langle i \rangle}$ 不受损。

$ $

\noindent \textbf{安全性：}固定实际秘密分布、误差分布、模数、维数和样本上限 $Q$。令 $\delta(T,Q)$ 表示时间为 $T$ 的区分器在真实 LWE 样本预言机与独立均匀二元组预言机之间的最大接受概率差。秘密重复使用，而每次查询的均匀掩码与误差均独立新鲜。这是判定假设，不仅是密钥恢复困难的断言。

在自适应左右加密游戏中，隐藏位 $\vartheta$ 从每次提交的消息对 $(m_0,m_1)$ 中选择一个消息。定义猜测偏差 $\epsilon=|\Pr[\widehat\vartheta=\vartheta]-1/2|$。归约自行选择 $\vartheta$，对每次查询从样本预言机取得 $(\vec a,z)$，返回 $(\vec a,z+\Delta m_\vartheta)$。真实样本下这恰好是加密；均匀样本下，即使给定此前全部记录，新鲜的 $z$ 加上任意消息偏移仍然均匀。因此逐步归纳可知，完整视图与 $\vartheta$ 独立。归约输出攻击者是否猜中，得到
\[
\epsilon\le\delta(T',Q),
\]
其中 $T'$ 包含模拟开销。由于假设中的优势已经允许 $Q$ 条样本，不需要再乘 $Q$。如果将 IND 优势定义成两个消息实验的输出概率差，对应上界是 $2\delta(T',Q)$。该结论涵盖经典加密接口中的重复和自适应选择明文，前提是不额外公开依赖秘密的材料。面对量子攻击者时，基础判定假设也必须涵盖量子区分器；此陈述不允许叠加态加密查询。

解密先减去依赖秘密的项，再通过舍入去除足够小的误差，除以 $\Delta$，最后模 $t$ 约减。噪声界与安全假设是两个独立要求：足够小而能正确解码，并不自动意味着足够大而能保证安全。

\subsection{RLWE 密码体制}
\label{subsec:lattice-scheme2}

RLWE 密码体制使用与 \tboxlabel{\ref*{subsec:lattice-scheme}} 相同的公式，但将 $\vec{s}, \vec{a}^{\langle i \rangle}, b^{\langle i \rangle}, m^{\langle i \rangle}, e^{\langle i \rangle}$ 替换为多项式 $S, A^{\langle i \rangle}, B^{\langle i \rangle}, M^{\langle i \rangle}, E^{\langle i \rangle}$，如下所示：

\begin{tcolorbox}[title={\textbf{\tboxlabel{\ref*{subsec:lattice-scheme2}} 基于格的 RLWE 密码体制}}]
\begin{itemize}
    \item \textbf{\underline{加密}：}$B^{\langle i \rangle} = S \cdot A^{\langle i \rangle} + \Delta \cdot M^{\langle i \rangle} + E^{\langle i \rangle}$，其中 $B^{\langle i \rangle}$ 和 $A^{\langle i \rangle}$ 对攻击者也是公开已知的，而 $S, M^{\langle i \rangle}, E^{\langle i \rangle}$ 则未知，仅秘密密钥持有者知晓。

    $ $
    
    \item \textbf{\underline{解密}：}$\left\lceil [B^{\langle i\rangle}-S\cdot A^{\langle i\rangle}]_q/\Delta\right\rfloor\bmod t=M^{\langle i\rangle}$ \\$\Big($前提是 $||E^{\langle i \rangle}||_{\infty} < \dfrac{\Delta}{2}$，即 $E^{\langle i \rangle}$ 的每个系数的绝对值都小于 $\dfrac{\Delta}{2}\Big)$

    $ $

    $\lceil  \rfloor_{\Delta}$ 表示对多项式每一项的系数分别进行舍入。

\end{itemize}



\end{tcolorbox}

上述安全归约在相应判定 RLWE 假设下适用于环元素对，也在实际 GLWE 假设下适用于 $k$ 个环掩码。将环乘法展开成系数方程，得到的是具有负循环结构的矩阵，并非独立均匀的 LWE 行。因此可以沿用证明模板，但标量 LWE 困难性定理并不会自动认证某个环或模结构参数族。
